\documentclass[12pt]{article}
\usepackage{amsmath}
\usepackage{geometry}
\geometry{a4paper, margin=1in}

\title{Lecture Scribe using Self-Consistency Prompting}
\author{}
\date{}

\begin{document}

\maketitle

\section*{Run 1}

\textbf{Reasoning:}

From the lecture images, the lecture starts with probability basics $\rightarrow$ definitions $\rightarrow$ axioms $\rightarrow$ propositions $\rightarrow$ assigning probability $\rightarrow$ joint probability $\rightarrow$ conditional probability $\rightarrow$ examples.

To generate the scribe, follow slide order and explain each concept step-by-step.

\subsection*{Scribe with reasoning}

\textbf{1. Lecture Topic Identification}

Topic: Joint Probability and Conditional Probability.

Lecture first reviews probability theory, then introduces joint and conditional probability.

\bigskip

\textbf{2. Basic Definitions reasoning}

Experiment = procedure producing result

Outcome = possible result

Event = set of outcomes

Sample space = all outcomes

\bigskip

\textbf{3. Axioms reasoning}

Probability rules:

\[
0 \le P(A) \le 1
\]

\[
P(S) = 1
\]

\[
P(A^c) = 1 - P(A)
\]

\[
P(A \cup B) = P(A) + P(B) - P(A \cap B)
\]

\bigskip

\textbf{4. Assigning probability reasoning}

Classical method → equally likely outcomes

Relative frequency → repeated experiment

\bigskip

\textbf{5. Joint probability reasoning}

Joint probability = probability of intersection

\[
P(A,B) = P(A \cap B)
\]

Example: card deck counting.

\bigskip

\textbf{6. Conditional probability reasoning}

\[
P(A|B) = \frac{P(A,B)}{P(B)}
\]

Product rule:

\[
P(A,B) = P(A|B)P(B)
\]

Used when events depend.

\bigskip

Conclusion: Lecture covers probability foundation → joint → conditional.


\newpage

\section*{Run 2}

\textbf{Reasoning:}

Follow slide order again.

definitions → axioms → propositions → assigning probability → joint → conditional → examples

\subsection*{Scribe with reasoning}

\textbf{1. Topic reasoning}

Lecture explains probability theory needed for joint and conditional probability.

\bigskip

\textbf{2. Definitions reasoning}

Experiment produces outcomes

Sample space contains all outcomes

Event is subset of sample space

\bigskip

\textbf{3. Axioms reasoning}

\[
0 \le P(A) \le 1
\]

\[
P(S) = 1
\]

Additivity for mutually exclusive events

\bigskip

\textbf{4. Propositions reasoning}

\[
P(A^c) = 1 - P(A)
\]

\[
P(A \cup B) = P(A) + P(B) - P(A \cap B)
\]

\bigskip

\textbf{5. Assign probability reasoning}

Classical counting

Relative frequency simulation

\bigskip

\textbf{6. Joint probability reasoning}

\[
P(A,B) = P(A \cap B)
\]

Example: cards.

\bigskip

\textbf{7. Conditional probability reasoning}

\[
P(A|B) = \frac{P(A,B)}{P(B)}
\]

Product rule:

\[
P(A,B) = P(A|B)P(B)
\]

Chain rule for many events.

\bigskip

Conclusion: Lecture goal = understand dependence between events.


\newpage

\section*{Run 3}

\textbf{Reasoning:}

Use same order again to ensure consistency.

\subsection*{Scribe with reasoning}

\textbf{1. Topic reasoning}

Lecture reviews probability then introduces joint and conditional probability.

\bigskip

\textbf{2. Definitions reasoning}

Sample space contains all outcomes

Events are subsets.

\bigskip

\textbf{3. Axioms reasoning}

Probability between 0 and 1

Whole space probability = 1

Additivity rule

\bigskip

\textbf{4. Propositions reasoning}

Complement rule

Union rule

General addition rule

\[
P(A^c)=1-P(A)
\]

\[
P(A \cup B)=P(A)+P(B)-P(A \cap B)
\]

\bigskip

\textbf{5. Assigning probability reasoning}

Classical method

Relative frequency method

\bigskip

\textbf{6. Joint probability reasoning}

\[
P(A,B)=P(A \cap B)
\]

Examples: cards, clothes.

\bigskip

\textbf{7. Conditional probability reasoning}

\[
P(A|B)=\frac{P(A,B)}{P(B)}
\]

Examples:

cards without replacement

poker

missing key

\bigskip

Conclusion: Lecture explains probability rules for dependent events.


\newpage

\section*{Selected Output}

\textbf{Selected Run: Run 2}

\bigskip

\textbf{Reason for selection}

\begin{itemize}
\item Same assumptions as other runs
\item Same slide order
\item Same formulas
\item Same reasoning steps
\item No missing concept
\item Highest agreement across runs
\end{itemize}

\end{document}